\documentclass[11pt]{article}
\usepackage[margin=1in]{geometry}
\usepackage{amsmath,amssymb,amsthm,mathtools}
\usepackage{booktabs}
\usepackage{graphicx}
\usepackage[hidelinks]{hyperref}
\usepackage{microtype}

\newtheorem{theorem}{Theorem}
\newtheorem{corollary}[theorem]{Corollary}
\newtheorem{lemma}[theorem]{Lemma}
\newcommand{\R}{\mathbb R}
\newcommand{\C}{\mathbb C}
\newcommand{\Sph}{\mathbb S^2}
\newcommand{\ip}[2]{\langle #1,#2\rangle}

\title{A two-channel instability region for the spinorial CKN inequality}
\author{Candidate partial result for arXiv:2504.16909}
\date{9 August 2026}

\begin{document}
\maketitle

\begin{center}
\fbox{\parbox{0.91\textwidth}{\textbf{Status.}
Candidate partial result; verifier verdict: \emph{likely valid}.  The theorem
below proves that mixing channels genuinely enlarges the instability region.
It does not determine the optimal region over all mixed channels.}}
\end{center}

\section{Source problem}

The source is J.~Dolbeault, M.~J.~Esteban, R.~L.~Frank, and M.~Loss,
\emph{The CKN inequality for spinors: symmetry and symmetry breaking},
arXiv:2504.16909v2 \cite{DEFL}.  In Remark~26 on PDF page~30, after deriving
single-channel instability regions, the authors write:

\begin{quote}
``It is not unlikely that by mixing several channels one might increase the
region of instability. Finding the optimal region is an open problem.''
\end{quote}

The complete source paragraph, the adjacent theorem, and the added-in-proof
novelty warning are reproduced below.

\begin{center}
\includegraphics[width=0.98\textwidth]{figures/open_problem_crop.png}
\end{center}

\section{Setup}

We use the source's logarithmic-cylinder formulation on
$\R\times\Sph$, where surface measure $d\omega$ on $\Sph$ is normalized to
be a probability measure.  For $\alpha>-1/2$, the symmetric critical point is
\[
 \phi_\star(s,\omega)=u_\star(s)\chi_0,
 \qquad \chi_0=\binom10.
\]
The source computes its real Hessian $\mathcal Q$.  We only need the following
specialization of that formula.  Let $\psi\in L^2(\Sph,\C^2)$ satisfy
\[
 \|\psi\|_{L^2(\Sph)}=1,
 \qquad \int_{\Sph}\ip{\chi_0}{\psi}\,d\omega=0,
\]
and let $w\in H^1(\R)$ be real.  Put
\[
 \delta(\psi):=\int_{\Sph}
 \left|\operatorname{Re}\ip{\chi_0}{\psi}\right|^2d\omega.
\]
Then the source Hessian gives
\begin{equation}\label{eq:reducedQ}
 \mathcal Q[w\psi]
 =\int_\R\left\{|w'|^2+\mu(\psi)|w|^2
 -A\bigl(1+(p-2)\delta(\psi)\bigr)
 \frac{|w|^2}{\cosh^2(Bs)}\right\}ds,
\end{equation}
where
\begin{equation}\label{eq:ABmu}
 A=\frac p2\left(\frac12-\alpha\right)^2,
 \quad
 B=\frac{p-2}{2}\left(\frac12-\alpha\right),
 \quad
 \mu(\psi)=\left\|\left(\sigma\!\cdot\!L+\frac12-\alpha\right)
 \psi\right\|_2^2.
\end{equation}
The zero-mean hypothesis makes the final positive rank-one term in the full
Hessian vanish.

The scalar minimization used repeatedly in the source is
\begin{equation}\label{eq:PT}
 \inf_{\|w\|_2=1}\int_\R
 \left(|w'|^2-C\frac{|w|^2}{\cosh^2(Bs)}\right)ds
 =-\frac14\left(\sqrt{4C+B^2}-B\right)^2.
\end{equation}
This is the ground-state formula for the P\"oschl--Teller operator.

\section{Proof intuition}

The two lowest channels used separately in the source have angular labels
$(k,m)=(1,1/2)$ and $(-2,1/2)$.  Their upper components are positive multiples
of the same real harmonic $Y_1^0$, while their lower components are opposite
multiples of $Y_1^1$.  There is therefore a unique normalized combination in
which the lower component cancels.  The result is the spinor
$(\sqrt{4\pi}Y_1^0,0)^T$.  Its real upper component carries the entire norm,
so the destabilizing nonlinear factor is maximal: $\delta=1$.  The cost of
this gain is only the weighted average of the two angular energies.  Once the
angular calculation is made, the radial minimization remains the same scalar
P\"oschl--Teller problem as in the source and yields an explicit threshold.

\clearpage
\section{The two-channel theorem}

\begin{theorem}[Explicit mixed-channel instability]\label{thm:main}
Let
\[
 -\frac12<\alpha<0,
 \qquad 2<p<6,
 \qquad d:=\frac12-\alpha.
\]
If
\begin{equation}\label{eq:threshold}
 p>p_{\mathrm{mix}}(\alpha)
 :=2\sqrt{1+\frac{2}{d^2}},
\end{equation}
then the symmetric critical point is linearly unstable.  Consequently,
\[
 C_{\alpha,p}<C^\star_{\alpha,p}.
\]
The interval $(p_{\mathrm{mix}}(\alpha),6)$ is nonempty for every
$\alpha\in(-1/2,0)$.  In particular, for
\[
 \frac{1-\sqrt3}{4}\leq\alpha<0
\]
this gives a nonempty instability region while Theorem~27 of the source gives
none from any single channel.  Hence channel mixing genuinely increases the
published instability region.
\end{theorem}

\begin{proof}
Write $Z_j:=\sqrt{4\pi}\,Y_1^j$.  Appendix~B of the source gives
\[
 \chi_1^{1/2}
 =\binom{\sqrt{2/3}\,Z_0}{\sqrt{1/3}\,Z_1},
 \qquad
 \chi_{-2}^{1/2}
 =\binom{\sqrt{1/3}\,Z_0}{-\sqrt{2/3}\,Z_1}.
\]
Define the normalized mixed channel
\begin{equation}\label{eq:psi}
 \psi:=\sqrt{\frac23}\,\chi_1^{1/2}
       +\sqrt{\frac13}\,\chi_{-2}^{1/2}.
\end{equation}
The displayed formulas show the exact cancellation
\[
 \psi=\binom{Z_0}{0}.
\]
Since $Y_1^0$ is real, normalized in ordinary surface measure, and has mean
zero, while the source uses probability surface measure, we have
\[
 \|Z_0\|_{L^2(d\omega)}=1,
 \quad \int_{\Sph}Z_0\,d\omega=0,
 \quad \delta(\psi)=\int_{\Sph}Z_0^2\,d\omega=1.
\]

The two constituent channels are orthonormal eigenvectors of
$\sigma\!\cdot\!L$, with eigenvalues $1$ and $-2$.  Therefore
\begin{align*}
 \mu(\psi)
 &=\frac23\left(\frac32-\alpha\right)^2
   +\frac13\left(-\frac32-\alpha\right)^2 \\
 &=\alpha^2-\alpha+\frac94
  =d^2+2.
\end{align*}
Insert these identities into \eqref{eq:reducedQ}.  The potential coefficient is
\[
 C=A(1+(p-2)\delta)=A(p-1)=\frac{p(p-1)}2d^2,
 \qquad B=\frac{p-2}{2}d.
\]
After optimizing $w$ by \eqref{eq:PT}, the least Hessian value in this
one-radial-profile, two-angular-channel subspace is
\begin{align}
 E_{\mathrm{mix}}(\alpha,p)
 &=d^2+2-\frac14\left(\sqrt{4C+B^2}-B\right)^2 \notag\\
 &=d^2+2-\frac{d^2}{16}
 \left(\sqrt{9p^2-12p+4}-(p-2)\right)^2.\label{eq:energy}
\end{align}

For $p>2$, the function
\[
 R(p):=\sqrt{9p^2-12p+4}-(p-2)
\]
is strictly increasing.  Indeed,
\[
 R'(p)=\frac{9p-6}{\sqrt{9p^2-12p+4}}-1>0,
\]
where the last inequality follows by squaring the positive quantities.
Equation \eqref{eq:energy} is zero exactly when
\[
 R(p)=4\sqrt{1+\frac2{d^2}}.
\]
Squaring once and simplifying gives the unique positive solution
\[
 p=2\sqrt{1+\frac2{d^2}}.
\]
Thus \eqref{eq:threshold} is equivalent to $E_{\mathrm{mix}}<0$.
The Hessian criterion from the source now implies linear instability and
$C_{\alpha,p}<C^\star_{\alpha,p}$.

Finally, $p_{\mathrm{mix}}(\alpha)<6$ is equivalent to $d>1/2$, which is
exactly $\alpha<0$.  The comparison with the source follows from its
Theorem~27, whose negative-$\alpha$ single-channel condition requires
$\alpha<(1-\sqrt3)/4$.
\end{proof}

\begin{corollary}[Reflected range]
Using the source symmetry $\alpha\mapsto-(1+\alpha)$, symmetric optimizers are
linearly unstable for
\[
 -1<\alpha<-\frac12,
 \qquad
 p>2\sqrt{1+\frac{2}{(\alpha+3/2)^2}}.
\]
\end{corollary}

\begin{proof}
Set $\widetilde\alpha:=-(1+\alpha)\in(-1/2,0)$ and apply
Theorem~\ref{thm:main}.  The source's reflection identity transfers the strict
inequality of constants back to $\alpha$, and
$1/2-\widetilde\alpha=\alpha+3/2$.
\end{proof}

\clearpage
\section{Why this is not yet the optimal region}

The construction above uses a common radial function in the two channels.
Allowing independent real profiles $w_+$ and $w_-$ produces the matrix
Schr\"odinger operator
\begin{equation}\label{eq:matrix}
 -\frac{d^2}{ds^2}I
 +\begin{pmatrix}
 (3/2-\alpha)^2&0\\0&(-3/2-\alpha)^2
 \end{pmatrix}
 -\frac{A}{\cosh^2(Bs)}
 \left[I+(p-2)vv^{\mathsf T}\right],
 \qquad
 v=\binom{\sqrt{2/3}}{\sqrt{1/3}}.
\end{equation}
For $\alpha\ne0$, the constant diagonal matrix and the rank-one matrix
$vv^{\mathsf T}$ do not commute.  A constant channel rotation therefore
cannot reduce \eqref{eq:matrix} to scalar P\"oschl--Teller operators.  The
unrestricted Hessian also contains analogous higher-angular blocks.  The
packet proves a complete theorem for the explicit subspace \eqref{eq:psi}, but
does not exclude lower eigenvalues in these larger blocks and hence does not
solve the source's optimal-region problem.

\section{Verification and numerical sanity check}

The companion script
\texttt{code/check\_two\_channel\_instability.py} checks:
\begin{itemize}
 \item the exact upper/lower component cancellation;
 \item the threshold identity at six values of $\alpha$;
 \item four points where the mixed energy is negative but both constituent
 one-channel energies are positive;
 \item one direct 200,001-point quadrature of the explicit P\"oschl--Teller
 ground-state Rayleigh quotient.
\end{itemize}
At $(\alpha,p)=(-0.10,5.50)$ the values are approximately
\[
 E_{\mathrm{mix}}=-0.3625,
 \quad E_{(1,1/2)}=0.69423,
 \quad E_{(-2,1/2)}=0.89738.
\]
The quadrature and closed mixed energy agree to better than $2\cdot10^{-9}$.
These computations are sanity checks and are not used in the proof.

\section{Novelty status and review recommendation}

A bounded search on 9 August 2026 covered the run's lightweight indexes and
parsed arXiv corpus, together with arXiv searches for the exact source title,
the authors Esteban and Frank, ``mixing several channels,'' and spinorial CKN
instability.  No separate public paper containing threshold
\eqref{eq:threshold} was found.  Novelty confidence is nevertheless only
moderate: version~2 of the source, posted 13 July 2026, adds a footnote saying
that recent work of Esteban and Frank suggests that mixing does increase the
region.  A human should check the two cancellation lines and the threshold
algebra, and should contact the source authors before any external originality
claim.

\begin{thebibliography}{9}
\bibitem{DEFL}
J.~Dolbeault, M.~J.~Esteban, R.~L.~Frank, and M.~Loss,
\emph{The CKN inequality for spinors: symmetry and symmetry breaking},
arXiv:2504.16909v2, 2026.
\end{thebibliography}

\end{document}
